\section{Introduction}\label{sec:introduction}

% Context
Software architecture defines the key design 
decisions that govern a system's structure and 
long-term evolution~\cite{bass2012,perry1992}.
As requirements change and delivery pressures increase, implementations
may drift from the intended architecture. In software engineering,
this divergence is known as architecture erosion.~\cite{desilva2012,perry1992}
This divergence manifests in layering violations, 
cyclic dependencies, and unauthorized couplings, and has been associated with declining
maintainability and increasing 
costs of change~\cite{passos2010,desilva2012}.

% --- 2. Problem / Motivation ---
One approach to mitigating architectural erosion is
to encode architectural constraints as executable 
checks that can be integrated into development
workflows and continuous integration (CI). 
In the evolutionary architecture
literature, these checks are often described 
as architectural \emph{fitness functions}, 
which automatically assess whether selected
structural properties are preserved as a system 
evolves~\cite{ford2017}.
Several tools support this approach across 
programming ecosystems, including
ArchUnit for JVM, import-linter for Python, 
and dependency-cruiser for
JavaScript/TypeScript.
These tools allow developers to specify constraints
on layering, dependencies, cycles, and visibility, 
and to configure checks
to fail a build when violations are detected. 
Integrating conformance
checks into CI has been proposed as a way to 
make architectural violations
a routine concern rather than an occasional 
discovery~\cite{pinto2017}.
However, prior research has reported that 
conformance tools remain
underused and that detected violations are rarely 
fixed~\cite{passos2010}.
This raises questions about how extensively 
executable architecture checks
are used in practice, which constraints they 
enforce, and whether they
remain active as software evolves.

% --- 3. Gap ---
Despite growing tooling support, empirical 
evidence on real-world
use and evolution of executable architecture 
rules remains limited.
Existing studies typically examine a small 
number of systems, focus on
a single ecosystem, or rely on handcrafted 
benchmarks~\cite{pinto2017,mujahid2024}.
Consequently, little is known about how 
widely these rules are adopted,
how they are integrated into development
workflows, and how their
configurations change over time through edits, 
exemptions, or disabling.
We define \emph{rule rot} as the exemption or disabling of an
executable architecture check, so that enforcement narrows or
ceases even when the rule specification remains.
Research on Architecture-as-Code (AaC) and 
design-rule mining has
also highlighted gaps in tooling and governance
~\cite{pontillo2026icsa,lucena2026ecsa}, while 
longitudinal studies of
open-source architectural governance have documented variation in
adoption and evidence of rule drift~\cite{larsen2026seaa}.
However, a systematic, repository-scale study spanning multiple
programming ecosystems is still needed to characterize adoption and
CI evidence, identify the architectural constraints encoded
in these rules, and examine how they evolve over time. How
source-level dependency and layering violation counts move from
before to after the first parseable ArchUnit rule, and what share
have zero evaluable violations one month before adoption versus
a pre-adoption violation that is absent at the adoption snapshot,
likewise remains undescribed. The first of those conditions does
not require a clean adoption snapshot.

% THIS PAPER.
To address these gaps, we conduct an empirical study
of 4,300 GitHub repositories containing artifacts
from ArchUnit, import-linter, and dependency-cruiser
across their respective programming
ecosystems. We extract executable architecture 
rules from these repositories, normalize them 
into a canonical schema, and reconstruct 
the histories of 11,098 rule files. Our findings 
reveal substantial differences 
across tools in both CI evidence and
rule content. We also observe considerable
variation in the 24-month cumulative incidence
of rule rot, while nearly half of the rule files
remain unchanged
throughout their observed histories. Finally, in the 317 ArchUnit repositories
with at least 12 months of history on each side of
the first parseable ArchUnit rule, we describe how
source-level dependency and layering violation counts
move from before that commit to after, and what share
have zero evaluable violations one month before adoption
versus a pre-adoption violation cleared at the adoption
snapshot. The first share does not require zero violations
at adoption.

% Research Question ----
Our study addresses four research questions covering 
four aspects: continuous integration (CI) evidence, 
rule content, rule evolution, and source-level violation 
trajectories associated with parseable ArchUnit rules.

\begin{itemize}
    \item \textbf{To what extent do repositories expose evidence of architecture-checker invocation in GitHub Actions workflows?}

    To answer RQ1, we hypothesize that fewer than
    half of import-linter repositories, and fewer than half of
    dependency-cruiser repositories, contain a GitHub Actions
    workflow that invokes the corresponding tool.

    \item \textbf{Which architectural constraints dominate
    executable rules, and how do they differ across tools?}

    To answer RQ2, we hypothesize that in each tool,
    dependency and layering constraints account for more than 70\%
    of classified constructs, excluding uncategorized constructs
    from the denominator.

    \item \textbf{How often are architecture rules exempted,
    disabled, or left unchanged?}

    To answer RQ3, we hypothesize that within 24 months, at least
    30\% of rule files experience a first observed exemption or
    disabling event.

    \item \textbf{How do dependency and layering violations
    change around ArchUnit adoption?}

    To address RQ4, we track changes in source-level
    dependency and layering violations before and
    after the introduction of the first parseable
    ArchUnit rule.
\end{itemize}

% --- 5. Contributions ---
In summary, we make the following contributions:
\begin{itemize}
  \item We construct and release ArchRules, a multi-ecosystem
  dataset comprising architecture rule files from 4,300 GitHub repositories,
  together with tools for canonical parsing and historical reconstruction
  of 11,098 rule files.

  \item We operationalize rule rot as the first observed
  exemption or disabling of an executable architecture check and
  characterize its occurrence over time. We also classify rules by
  architectural constraint to support systematic comparisons across tools.

  \item We compare CI evidence and rule content across tools.
  Tool invocation in a GitHub Actions workflow is measured for
  import-linter and dependency-cruiser. For ArchUnit, the evidence
  is a Maven or Gradle build or test proxy and does not establish
  that ArchUnit ran. We also document the architectural constraints
  each tool encodes.

  \item We describe how source-level dependency and layering
  violation counts move from before to after the first parseable
  ArchUnit rule in the 317 repositories with at least 12 months
  of history on each side. We report the share with zero
  evaluable violations one month before adoption, the share
  that clear a pre-adoption violation at the adoption snapshot,
  and the share with violations at both snapshots.
\end{itemize}